\section{The finite carrier reduction and its two rational cones}
\label{sec:carrier-cones}
The finite local input consists of the eight master groups and all their
literal normal kernels. The respective numbers of kernels are
$27,13,28,2911,2911,2911,2292,2915$. Their generators and all
central-involution extension edges are supplied in the carrier normal
certificate. Lemma~\ref{lem:binary-coverage} proves completeness from
these edges. For each kernel the checked tuple is
$(k,n,m,a_2,c,g)$ as in Lemma~\ref{lem:carrier-transition} and
\eqref{eq:carrier-epi}. In particular the maximum defining $m$ is
computed over the same complete normal list, and the original quotient
is used for $c,g$.

For the degree-sixteen groups the identities $\Phi(G)=G'$ permit
$a_2\leq m$: here $N^2[N,G]\leq N\cap G'$. Their rows are dominated
by the following twelve scale-two rows; the other three blocks below
are the scale-one rows of $X,J,P$. Domination is simultaneous on the
pair supports and both marks $k,q=\max(m,a_2)$, including every cross
column between degrees eight and sixteen.
\[
\begin{array}{c|rrrrrr}
H16&k&n&m&a_2&c&g\\\hline
0&0&0&0&0&1&6\\
1&1&1&1&1&2&5\\
2&1&2&1&1&4&4\\
3&2&3&2&2&3&3\\
4&2&4&2&2&4&2\\
5&3&6&3&3&6&0\\
6&4&6&4&4&5&0\\
7&4&7&4&4&4&0\\
8&4&8&4&4&3&0\\
9&5&11&3&3&1&0\\
10&5&11&4&4&0&0\\
11&6&12&3&3&0&0\\
\end{array}
\]
\[
\begin{array}{c|rrrrrr}
X&k&n&m&a_2&c&g\\\hline
0&0&0&0&0&1&3\\
1&1&1&1&0&2&2\\
2&1&2&1&1&2&1\\
3&1&3&1&1&1&1\\
4&1&4&2&2&2&0\\
5&2&3&2&1&3&0\\
6&2&4&2&1&2&0\\
7&2&5&2&2&1&0\\
8&3&6&2&2&0&0\\
\end{array}
\]
\[
\begin{array}{c|rrrrrr}
J&k&n&m&a_2&c&g\\\hline
0&0&0&0&0&1&3\\
1&1&1&1&0&1&2\\
2&1&2&1&1&2&1\\
3&1&3&1&1&1&1\\
4&1&3&1&1&3&0\\
5&1&4&1&1&2&0\\
6&1&5&1&2&1&0\\
7&2&4&1&1&2&0\\
8&2&5&1&1&1&0\\
9&2&6&1&2&0&0\\
\end{array}
\]
\[
\begin{array}{c|rrrrrr}
P&k&n&m&a_2&c&g\\\hline
0&0&0&0&0&1&4\\
1&1&1&1&0&1&3\\
2&1&2&1&1&2&2\\
3&1&3&1&1&2&1\\
4&1&4&1&1&1&1\\
5&1&5&2&2&2&0\\
6&2&4&2&1&3&0\\
7&2&5&2&1&2&0\\
8&2&6&2&2&1&0\\
9&3&7&2&2&0&0\\
\end{array}
\]

The normalized representatives are
\[
\begin{array}{c|rrrrrr|c}
 &k&n&m&a_2&c&g&s\\\hline
 H_9&5&11&3&3&1&0&2\\
 X_{\rm whole}&3&6&2&2&0&0&1\\
 P_0&0&0&0&0&1&4&1\\
 P_1&1&1&1&0&1&3&1\\
 P_2&1&2&1&1&2&2&1\\
 P_6&2&4&2&1&3&0&1\\
 P_{\rm whole}&3&7&2&2&0&0&1
\end{array}
\]
For completeness, the following list gives a dominating representative
for each row in each displayed block, in its displayed order.
The last representative has exactly the same normalized pair row and
marks as $X_{\rm whole}$, so they are merged. All comparisons here are
finite rational inequalities of the explicit support functions.
\[
H16:\quad (P_0,P_1,P_2,P_2,P_2,P_6,P_6,P_6,P_6,H_9,X_{\rm whole},X_{\rm whole}).
\]
\[
X:\quad (P_0,P_2,P_2,P_2,P_6,P_6,P_6,P_6,X_{\rm whole}).
\]
\[
J:\quad (P_0,P_1,P_2,P_2,P_2,P_2,P_6,P_6,P_6,X_{\rm whole}).
\]
\[
P:\quad (P_0,P_1,P_2,P_2,P_2,P_6,P_6,P_6,P_6,X_{\rm whole}).
\]

The reduction of the actual degree-sixteen tuples to the twelve rows
also checks all degree-eight columns, and the two marks, in the
committed transition certificate. Thus the preceding reduction does
not presume that domination inside a smaller submatrix survives the
addition of new colours. Scaling and merging give exactly
\eqref{eq:carrier-B}.

It remains to prove the scalar inequality used in
Theorem~\ref{thm:carrier-cross}. Put $p=\sum_i p_i$ and
\[
 \mathcal D=R^2/4+(139/328)Rp+(627/2624)p^2-\mathcal E.
\]
In the first cone let $(L,J,S)=(\ell,j-2\ell,R-j)$, and in the second
let $(A,C,S)=(j,2\ell-j,R-2\ell)$. These are nonnegative in their
respective cones. Let $z$ consist of $p_1,\ldots,p_6$ followed by the three
corresponding cone coordinates.
Expansion gives $5248\mathcal D=z^{\mathsf T}M_\nu z$.
The upper left $6\times6$ block, common to both $M_\nu$, is
\[
M_0=\begin{pmatrix}
1049&1008&-140&106&-58&24\\
1008&1254&-550&-222&-386&-222\\
-140&-550&1254&598&434&-386\\
106&-222&598&762&598&-58\\
-58&-386&434&598&598&-58\\
24&-222&-386&-58&-58&270
\end{pmatrix}.
\]
Write $M_\nu=\begin{pmatrix}M_0&D_\nu\\D_\nu^{\mathsf T}&E_\nu\end{pmatrix}$.
The remaining blocks are
\[
D_A=\begin{pmatrix}
-400&-528&1112\\
-1056&-856&1112\\
2224&1112&1112\\
912&456&1112\\
912&456&1112\\
-400&-200&1112
\end{pmatrix},\qquad E_A=\begin{pmatrix}
5248&1312&-1312\\
1312&1312&-1312\\
-1312&-1312&1312
\end{pmatrix}.
\]
\[
D_B=\begin{pmatrix}
-200&-200&1112\\
-528&-528&1112\\
1112&1112&1112\\
456&456&1112\\
456&456&1112\\
-200&-200&1112
\end{pmatrix},\qquad E_B=\begin{pmatrix}
1312&0&-656\\
0&3936&1968\\
-656&1968&1312
\end{pmatrix}.
\]
The following rational matrices give positive-semidefinite parts.
\begingroup\setlength\arraycolsep{2pt}\scriptsize
\[
P_A=\begin{pmatrix}
817&804&-548&-328&-393&-\frac{540}{11}&-896&-688&688\\
804&1254&-550&-288&-386&-222&-1129&-856&856\\
-548&-550&1254&423&434&-386&1836&1112&-1112\\
-328&-288&423&577&237&-\frac{3987}{44}&563&323&-323\\
-393&-386&434&237&490&-58&646&363&-363\\
-\frac{540}{11}&-222&-386&-\frac{3987}{44}&-58&270&-\frac{19363}{44}&-200&200\\
-896&-1129&1836&563&646&-\frac{19363}{44}&4645&1312&-1312\\
-688&-856&1112&323&363&-200&1312&1312&-1312\\
688&856&-1112&-323&-363&200&-1312&-1312&1312
\end{pmatrix}.
\]
\[
P_B=\begin{pmatrix}
729&721&-469&-236&-318&-\frac{602}{11}&-416&-480&86\\
721&1254&-550&-296&-386&-222&-565&-585&-2\\
-469&-550&1254&433&434&-386&917&950&-407\\
-236&-296&433&542&241&-\frac{4061}{44}&275&197&-56\\
-318&-386&434&241&469&-58&303&231&-65\\
-\frac{602}{11}&-222&-386&-\frac{4061}{44}&-58&270&-\frac{2413}{11}&-\frac{10005}{44}&\frac{8581}{44}\\
-416&-565&917&275&303&-\frac{2413}{11}&1254&-134&-813\\
-480&-585&950&197&231&-\frac{10005}{44}&-134&2690&591\\
86&-2&-407&-56&-65&\frac{8581}{44}&-813&591&891
\end{pmatrix}.
\]
\endgroup
Subtraction gives $M_A-P_A\geq0$ entrywise, with entries at most
$2224$, and $M_B-P_B\geq0$ entrywise, with entries at most $1519$.
These assertions are direct comparisons of the displayed rational
entries. Both matrices annihilate
$v=(0,17,21,0,0,44,0,0,0)$; $P_A$ also annihilates
$w=(0,0,0,0,0,0,0,1,1)$.
Delete coordinates six and nine from $P_A$. The successive leading
principal determinants of the retained matrix are
\[
\begin{gathered}
817,\quad378102,\quad335067392,\quad132631407594,\\
35034163435774,\quad63581756789450739,\quad2100371145348629273.
\end{gathered}
\]
They are positive. The omitted coordinates of $(v,w)$ form
$\operatorname{diag}(44,1)$, so subtracting their unique linear
combination carries any vector into this retained subspace without
changing its quadratic value. Thus $P_A$ is positive semidefinite.
Delete coordinate six from $P_B$. Its retained leading principal
determinants are
\[
\begin{gathered}
729,\quad394325,\quad370093856,\quad140632346331,\\
36346448139435,\quad19760569015182064,\\
7216647669027584752,\quad1335792379687239644084.
\end{gathered}
\]
Again they are positive, and the omitted coordinate of $v$ is $44$,
so the same kernel elimination proves that $P_B$ is positive
semidefinite. Hence $z^{\mathsf T}M_\nu z\geq0$ for every nonnegative
$z$. The two cones exhaust all $0\leq j\leq R$, $0\leq\ell\leq R/2$,
including $j<\ell$, proving the required scalar inequality.
